% BANKS-SOURCE: pdf=216 print=206 kind=chapter id=chapter10
\BanksChapter{10}{Instantons and solitons}

% BANKS-SOURCE: pdf=216 print=206 kind=prose
Much of this book has been devoted to perturbation theory around Gaussian fixed
points of the renormalization group. Such a perturbation theory can always be
reorganized into a semi-classical approximation to the functional integral. The
semi-classical expansion can also give us evidence about non-perturbative effects
in quantum field theory, and this chapter is an introduction to the relevant
technology. It is a very brief introduction to a very large subject. Readers who
wish to pursue the subject of instantons and solitons in depth should consult
\cite{banks-ref-135,banks-ref-136,banks-ref-137,banks-ref-138}.

% BANKS-SOURCE: pdf=216 print=206 kind=section id=10.1
\subsection{The most probable escape path}

% BANKS-SOURCE: pdf=216 print=206 kind=prose
We begin by recalling the WKB approximation to the problem of quantum tunneling
in a system with $N$ degrees of freedom, $q^i$.

The Schrödinger equation is

% BANKS-SOURCE: pdf=216 print=206 kind=equation id=10.1
\begin{equation}
  \left[-\frac{g^2}{2}\frac{\partial^2}{\partial(q^i)^2}
    +\frac{1}{g^2}V(q)\right]\psi=E\psi.
  \label{eq:10.1}
\end{equation}

% BANKS-SOURCE: pdf=216 print=206 kind=prose
For simplicity we have made all the variables, including the energy,
dimensionless, set $\hbar=1$, and introduced a dimensionless parameter $g^2$.
The WKB approximation is valid when $g^2$ is small. The reader should verify
that this Schrödinger equation would be appropriate for a system whose classical
action is

% BANKS-SOURCE: pdf=216 print=206 kind=equation id=10.2
\begin{equation}
  S=\int\!\dd t\left[\frac{(\dot q^i)^2}{2g^2}-\frac{V}{g^2}\right].
  \label{eq:10.2}
\end{equation}
\BanksImplicitHook{I-CH10-001}

% BANKS-SOURCE: pdf=216 print=206 kind=prose
For small $g^2$, an approximate solution to the Schrödinger equation is

% BANKS-SOURCE: pdf=216 print=206 kind=equation id=10.3
\begin{equation}
  \psi=\ee^{\frac{\ii}{g^2}S},
  \label{eq:10.3}
\end{equation}

where

% BANKS-SOURCE: pdf=216 print=206 kind=equation id=10.4
\begin{equation}
  \left(\frac{\partial S}{\partial q^i}\right)^2=2(E-V).
  \label{eq:10.4}
\end{equation}

% BANKS-SOURCE: pdf=216 print=206 kind=prose
The last equation is the Hamilton--Jacobi equation for the classical system with
action (10.2). The integral curves of this first-order system are solutions of the
classical equations of motion, and

% BANKS-SOURCE: pdf=216 print=206 kind=equation id=10.5
\begin{equation}
  \frac{\partial S}{\partial q^i}=p_i=\dot q_i.
  \label{eq:10.5}
\end{equation}

% BANKS-SOURCE: pdf=217 print=207 kind=prose
Here $q^i(t)$ is the classical trajectory that goes through the point $q^i$ at
time $t$. The velocities are real only if $E>V$, but we can use this
approximation to solve the Schrödinger equation in the tunneling region where
$V>E$. Since the coordinates are still real, this corresponds to continuing to
imaginary time.

The imaginary-time equations of motion are perfectly good differential equations.
In fact, they correspond to Newtonian motion in the potential $-V$. The relation
between solutions to this equation and the WKB approximation to the Schrödinger
equation is the same as that between ordinary classical trajectories and the
Hamilton--Jacobi equation. In particular, we have the usual result that, at any
point in the tunneling region,

% BANKS-SOURCE: pdf=217 print=207 kind=equation id=10.6
\begin{equation}
  S=\pm\ii\int_{q_0}^{q}p_i\,\dd q^i,
  \label{eq:10.6}
\end{equation}

% BANKS-SOURCE: pdf=217 print=207 kind=prose
where the integral is taken over an imaginary-time classical trajectory that goes
through the indicated points. This is just the Euclidean action of the trajectory.
We generally take $q_0$ to be the turning point of the trajectory, where it hits
the boundary at which real classical motion can begin. The two signs correspond
to linearly independent solutions of the Schrödinger equation, one of which is
exponentially falling and the other exponentially increasing, as we penetrate
into the tunneling region. Boundary conditions on the other side of the
tunneling barrier (or at infinity if it is not a finite barrier, but an infinitely
rising wall) fix the coefficient of this growing solution to be small, so that the
two are of the same order of magnitude in the middle of the tunneling region.

For small $g$, the wave function is exponentially small in the tunneling region,
and it will be maximized in the vicinity of the path of minimal Euclidean action
which traverses the barrier. In minimizing the action we have to search over all
solutions that traverse the barrier between two classically allowed regions in
configuration space. This includes a search over all end points of the solution
after tunneling. However, once we have found the minimum, there is another
solution in which we first follow the most probable escape path, and then retrace
it to its origin. This solution, called the *bounce*, minimizes the action subject
to the constraint that it starts and ends at the initial stationary point of the
potential. There is no need to search over possible turning points in the second
classically allowed region. The action of the bounce solution directly gives the
logarithm of the probability for tunneling through the barrier.

% BANKS-SOURCE: pdf=217 print=207 kind=section id=10.2
\subsection{Instantons in quantum mechanics}

% BANKS-SOURCE: pdf=217 print=207 kind=prose
Consider one-dimensional quantum mechanics with a potential $V(x)$ that has two
or more minima. We will concentrate on three cases: two degenerate minima, a
periodic potential with an infinite number of degenerate minima, and two
non-degenerate minima. The Euclidean (imaginary-time) action has the form

% BANKS-SOURCE: pdf=217 print=207 kind=display
\[
  \frac{1}{g^2}\int\!\dd t\left[\frac{1}{2}
  \left(\frac{\dd x}{\dd t}\right)^2+V(x)\right].
\]
